\needspace{7\baselineskip}
\section{Fundamentos de Estatística}
\label{sec:estatistica}

Estatística ajuda a resumir, comparar e interpretar dados. Antes de calcular uma medida, identifique a unidade de observação, as variáveis e a pergunta que o resultado deve responder.

\subsection{População, amostra e unidade de observação}

A \textbf{população} reúne as unidades sobre as quais queremos obter informação. Uma \textbf{amostra} é um subconjunto observado dessa população. Se a pergunta for sobre todas as pessoas inscritas em um curso, os resultados de um questionário respondido por apenas parte delas descrevem os respondentes. A generalização depende de como a amostra foi obtida e de quem ficou de fora.

Uma linha pode representar uma pessoa, uma venda ou uma medição. Se uma pessoa aparece em várias linhas, o número de linhas e o número de pessoas respondem a perguntas diferentes. Na base de notas deste guia, cada linha corresponde a um participante fictício.

\subsection{Tipos de variável}

Uma variável \textbf{categórica} identifica grupos ou níveis, como o curso de um participante. Uma variável \textbf{numérica} representa uma contagem ou medida, como o número de atividades entregues ou as horas de estudo. Um identificador numérico, como o código de um participante, serve para distinguir registros. Calcular sua média não responde a uma pergunta sobre a turma.

Uma variável categórica pode ser \textbf{nominal}, como o curso, ou \textbf{ordinal}, como níveis baixo, médio e alto. Entre variáveis numéricas, uma contagem é \textbf{discreta}, enquanto uma medida como tempo é geralmente tratada como \textbf{contínua}, mesmo que os valores tenham sido registrados com arredondamento.

\subsection{Frequências absolutas e relativas}

A frequência absoluta é a contagem de observações em uma categoria. A relativa divide essa contagem pelo total considerado. Neste exemplo, duas das cinco observações pertencem à categoria \codigo{"Manha"}, portanto sua frequência relativa é \(2/5=0{,}4\), ou \(40\%\).

\begin{exemplo}
from collections import Counter

turnos = ["Manha", "Tarde", "Manha", "Noite", "Tarde"]
contagens = Counter(turnos)

for turno, total in contagens.items():
    proporcao = total / len(turnos)
    print(f"{turno}: {total} ({proporcao:.0%})")
# Manha: 2 (40%), Tarde: 2 (40%), Noite: 1 (20%)
\end{exemplo}

Uma tabela de frequências facilita a comparação entre categorias. Se houver dados ausentes, declare o denominador usado. Por exemplo, dividir pelo total de respostas válidas descreve a distribuição entre quem respondeu.

\subsection{Medidas de tendência central}

Para \(n\) observações, a média aritmética é \(\bar{x}=\frac{1}{n}\sum_{i=1}^n x_i\). A mediana é o valor central dos dados ordenados. Se \(n\) for par, ela é a média dos dois valores centrais. A moda é o valor mais frequente. Um conjunto pode ter mais de uma moda quando há empate na maior frequência.

\begin{exemplo}
from collections import Counter
from statistics import median

valores = [10, 12, 15, 20, 30]
print(sum(valores) / len(valores))  # 17.4
print(median(valores))             # 15

categorias = ["A", "A", "B", "C", "A", "B"]
moda, frequencia = Counter(categorias).most_common(1)[0]
print(moda, frequencia)            # A 3
\end{exemplo}

Se o último valor passar de \(30\) para \(100\), a média passa de \(17{,}4\) para \(31{,}4\) e a mediana continua em \(15\). O valor extremo pesa no cálculo da média, enquanto a posição central permanece igual.


Se houver empate na maior frequência, \codigo{statistics.multimode()} retorna todos os valores empatados.

\begin{exemplo}
from statistics import multimode

print(multimode([6, 6, 8, 8, 9]))  # [6, 8]
\end{exemplo}

\subsection{Medidas de dispersão}

A \textbf{amplitude} é \(x_{\max}-x_{\min}\). Para os cinco valores acima, ela é \(20\). Se \(30\) for substituído por \(100\), ela passa a \(90\).

O \textbf{intervalo interquartil} é \(\mathrm{IQR}=Q_3-Q_1\), em que \(Q_1\) e \(Q_3\) são os percentis 25 e 75. Ele mede a largura da região central da distribuição. Os limites \(Q_1-1{,}5\,\mathrm{IQR}\) e \(Q_3+1{,}5\,\mathrm{IQR}\) ajudam a sinalizar observações para investigação. Registre a convenção usada para calcular quartis, pois métodos de interpolação podem produzir resultados diferentes.


Para os dados \codigo{[10, 12, 15, 20, 30]}, o método linear produz \(Q_1=12\), \(Q_3=20\) e \(\mathrm{IQR}=8\). O limite superior fica em \(32\). Se o \(30\) for substituído por \(100\), os quartis continuam iguais neste exemplo, e \(100\) fica acima do limite superior. A regra indica um ponto a investigar. Seu contexto determina como interpretá-lo.

\needspace{6\baselineskip}
Para descrever uma população finita de \(N\) valores com média \(\mu\), a variância é
\[
  \sigma^2 = \frac{1}{N}\sum_{i=1}^{N}(x_i-\mu)^2.
\]
Para a variância amostral corrigida, com \(n>1\) observações, usamos
\[
  s^2 = \frac{1}{n-1}\sum_{i=1}^{n}(x_i-\bar{x})^2.
\]
O desvio padrão é a raiz quadrada da variância e tem a mesma unidade dos dados originais.

\begin{exemplo}
import numpy as np

valores = np.array([10, 12, 15, 20, 30])
q1, q3 = np.percentile(valores, [25, 75], method="linear")
print(q1, q3, q3 - q1)  # 12.0 20.0 8.0

print(round(np.var(valores, ddof=0), 2))  # 51.04
print(round(np.var(valores, ddof=1), 2))  # 63.8
print(round(np.std(valores, ddof=0), 2))  # 7.14
print(round(np.std(valores, ddof=1), 2))  # 7.99
\end{exemplo}

\nota{Atenção ao denominador}{O NumPy usa \codigo{ddof=0} por padrão. Já \codigo{Series.std()} do pandas usa \codigo{ddof=1}. Ao comparar resultados, fixe o mesmo \codigo{ddof} e aplique o mesmo tratamento a valores ausentes \cite{numpy,pandas}.}

\subsection{Correlação de Pearson}

O coeficiente de Pearson mede a direção e a intensidade da associação \textbf{linear} entre observações pareadas de duas variáveis numéricas. Se ambas tiverem variância positiva, seu valor fica entre \(-1\) e \(1\):
\[
r=
\frac{\sum_{i=1}^{n}(x_i-\bar{x})(y_i-\bar{y})}
{\sqrt{\sum_{i=1}^{n}(x_i-\bar{x})^2}\sqrt{\sum_{i=1}^{n}(y_i-\bar{y})^2}}.
\]

\begin{exemplo}
import numpy as np

x = np.array([1, 2, 3, 4, 5])
y = np.array([2, 4, 5, 4, 6])
print(round(np.corrcoef(x, y)[0, 1], 2))  # 0.85
\end{exemplo}

Valores próximos de zero indicam pouca associação linear. Um gráfico de dispersão ajuda a identificar relações curvas, agrupamentos e pontos influentes. Uma associação observada, mesmo forte, é insuficiente para estabelecer causalidade.

\subsection{Probabilidade (aprofundamento)}

\paragraph{Distribuição binomial.}
Se \(X\) conta sucessos em \(n\) tentativas independentes, com a mesma probabilidade de sucesso \(p\) em cada uma, então
\[
  P(X=k)=\binom{n}{k}p^k(1-p)^{n-k},
  \qquad k=0,1,\ldots,n.
\]

\paragraph{Distribuição de Poisson.}
A distribuição de Poisson é usada para modelar contagens em um intervalo quando o processo pode ser aproximado por uma taxa constante e contagens independentes em intervalos disjuntos. Se \(\lambda>0\) é o número esperado de ocorrências no intervalo, então
\[
  P(X=k)=\frac{e^{-\lambda}\lambda^k}{k!},
  \qquad k=0,1,2,\ldots.
\]

Nos exemplos abaixo, calculamos a probabilidade de exatamente dois sucessos em dez tentativas com \(p=0{,}3\) e a probabilidade de exatamente duas ocorrências em um intervalo cuja contagem esperada é três \cite{scipy}.

\begin{exemplo}
from scipy.stats import binom, poisson

print(round(binom.pmf(2, n=10, p=0.3), 4))  # 0.2335
print(round(poisson.pmf(2, mu=3), 4))      # 0.2240
\end{exemplo}
